\documentclass[11pt,a4paper]{article}
\usepackage[utf8]{inputenc}
\usepackage[vietnamese]{babel}
\usepackage{amsmath,amssymb}
\usepackage{geometry}
\geometry{a4paper, total={180mm,260mm}, left=15mm, right=15mm, top=18mm, bottom=18mm}
\usepackage{tikz}
\usetikzlibrary{shapes.geometric, arrows.meta, positioning, calc}
\usepackage{tcolorbox}
\tcbuselibrary{skins,breakable}
\usepackage{xcolor}
\usepackage{multicol}
\usepackage{array}
\usepackage{fontawesome}
\usepackage{fancyhdr}

\pagestyle{fancy}
\fancyhf{}
\renewcommand{\headrulewidth}{0.5pt}
\renewcommand{\footrulewidth}{0.5pt}
\fancyhead[L]{\small\bfseries\color{teal}\faBook\ SỔ TAY KIẾN THỨC TOÁN 3}
\fancyhead[R]{\small\bfseries\color{teal}Lớp Toán Thầy Lĩnh -- loptoanthaylinh.vn}
\fancyfoot[C]{\small\color{gray}-- Sổ tay kiến thức: Bài 63. Phép cộng trong phạm vi 100 000 --}

\definecolor{primary}{HTML}{005F73}
\definecolor{secondary}{HTML}{0A9396}
\definecolor{accent}{HTML}{EE9B00}
\definecolor{danger}{HTML}{AE2012}
\definecolor{boxbg}{HTML}{F4F6F8}

\begin{document}

% --- BANNER TIÊU ĐỀ ---
\begin{tcolorbox}[
  enhanced,
  colback=white,
  colframe=primary,
  arc=3mm,
  boxrule=1.2pt,
  left=4mm, right=4mm, top=2.5mm, bottom=2.5mm,
  drop shadow=black!10!white
]
  \centering
  {\Large\bfseries\color{primary}\faBookmark\ BÀI 63: PHÉP CỘNG TRONG PHẠM VI 100 000}\\[1.5mm]
  {\small\itshape\color{black!75}Kĩ thuật đặt tính 5 chữ số thẳng cột, cộng từ phải sang trái, cộng có nhớ và tính nhẩm số tròn chục nghìn trên 1 trang A4}
\end{tcolorbox}

\vspace{-2mm}


\noindent
\begin{minipage}[t]{0.49\textwidth}\vspace{0pt}
\begin{tcolorbox}[
  enhanced,
  colback=blue!4!white,
  colframe=blue!75!black,
  arc=2mm, boxrule=0.8pt,
  left=2mm, right=2mm, top=1.8mm, bottom=1.8mm,
  title={\footnotesize\bfseries\color{white}\faKey\ 1. Quy trình Đặt tính rồi tính}
]
{\footnotesize
\begin{itemize}\setlength{\itemsep}{1.5pt}\setlength{\leftskip}{-2.5mm}
  \item \textbf{Đặt tính:} Viết các số hạng sao cho các chữ số ở cùng một hàng thẳng cột với nhau:
  $$\text{ĐV thẳng ĐV, Chục thẳng Chục, \dots, Chục nghìn thẳng Chục nghìn}$$
  \item \textbf{Thực hiện phép cộng:} Cộng lần lượt từng hàng từ \textbf{phải sang trái}:
  $$\text{Đơn vị} \to \text{Chục} \to \text{Trăm} \to \text{Nghìn} \to \text{Chục nghìn}$$
  \item \textbf{Cộng có nhớ:} Nếu tổng của một hàng $\ge 10$, viết chữ số đơn vị, nhớ $1$ sang hàng tiếp theo bên trái.
\end{itemize}
}
\end{tcolorbox}
\end{minipage}
\hfill
\begin{minipage}[t]{0.49\textwidth}\vspace{0pt}
\begin{tcolorbox}[
  enhanced,
  colback=teal!4!white,
  colframe=teal!75!black,
  arc=2mm, boxrule=0.8pt,
  left=2mm, right=2mm, top=1.8mm, bottom=1.8mm,
  title={\footnotesize\bfseries\color{white}\faCalculator\ 2. Tính nhẩm các số tròn chục nghìn}
]
{\footnotesize
\begin{itemize}\setlength{\itemsep}{1.5pt}\setlength{\leftskip}{-2.5mm}
  \item \textbf{Cộng tròn chục nghìn:}
  $$30\,000 + 40\,000 = \mathbf{70\,000}$$
  (Nhẩm: $3\text{ chục nghìn} + 4\text{ chục nghìn} = 7\text{ chục nghìn}$).
  $$50\,000 + 50\,000 = \mathbf{100\,000}$$
  \item \textbf{Cộng tròn nghìn:}
  $$42\,000 + 5000 = \mathbf{47\,000}$$
\end{itemize}
}
\end{tcolorbox}
\end{minipage}

\vspace{-2.5mm}
% --- KHỐI 3: SƠ ĐỒ TƯ DUY TIKZ ---
\noindent
\begin{tcolorbox}[
  enhanced,
  colback=orange!4!white,
  colframe=orange!80!black,
  arc=2mm, boxrule=0.8pt,
  left=2mm, right=2mm, top=1.5mm, bottom=1.5mm,
  title={\footnotesize\bfseries\color{white}\faSitemap\ SƠ ĐỒ TƯ DUY: TIẾN TRÌNH CỘNG 5 CHỮ SỐ}
]
\centering
\begin{tikzpicture}[
  box/.style={draw, rounded corners=2pt, align=center, font=\scriptsize\bfseries, inner sep=2.5pt, minimum height=0.75cm},
  arr/.style={-{Stealth[length=2mm]}, thick}
]
  \node[box, fill=purple!15, draw=purple!80!black, text width=2.1cm] (p1) at (0, 0) {ĐƠN VỊ\\Cộng trước\\Nhớ sang chục};
  \node[box, fill=teal!15, draw=teal!80!black, text width=2.1cm] (p2) at (2.8, 0) {HÀNG CHỤC\\Cộng thêm nhớ\\Nhớ sang trăm};
  \node[box, fill=blue!12, draw=blue!80!black, text width=2.1cm] (p3) at (5.6, 0) {HÀNG TRĂM\\Cộng thêm nhớ\\Nhớ sang nghìn};
  \node[box, fill=green!15, draw=green!70!black] (p4) at (8.4, 0) {HÀNG NGHÌN\\Cộng thêm nhớ\\Nhớ sang chục nghìn};
  \node[box, fill=orange!15, draw=orange!80!black, text width=2.3cm] (p5) at (11.2, 0) {CHỤC NGHÌN\\Cộng nhớ cuối\\Viết kết quả};

  \draw[arr, purple!80!black] (p1) -- (p2);
  \draw[arr, teal!80!black] (p2) -- (p3);
  \draw[arr, blue!80!black] (p3) -- (p4);
  \draw[arr, green!70!black] (p4) -- (p5);
\end{tikzpicture}
\end{tcolorbox}

\vspace{-2.5mm}
% --- KHỐI 4: VÍ DỤ & MẸO TRÁNH BẪY ---
\noindent
\begin{minipage}[t]{0.49\textwidth}\vspace{0pt}
\begin{tcolorbox}[
  enhanced,
  colback=green!4!white,
  colframe=green!60!black,
  arc=2mm, boxrule=0.8pt,
  left=2mm, right=2mm, top=1.8mm, bottom=1.8mm,
  title={\footnotesize\bfseries\color{white}\faCheckCircle\ Ví dụ thực hành chuẩn mực}
]
{\footnotesize
\textbf{1. Đặt tính cộng có nhớ:}\\
$\begin{array}{r@{}l}
  48256 & \\
  +\quad 27185 & \\
  \hline
  \mathbf{75441} &
\end{array}$
\qquad
\begin{minipage}[c]{0.5\linewidth}
-- $6 + 5 = 11$, viết $1$ nhớ $1$.\\
-- $5 + 8 = 13$, thêm $1$ là $14$, viết $4$ nhớ $1$.\\
-- $2 + 1 = 3$, thêm $1$ là $4$, viết $4$.\\
-- $8 + 7 = 15$, viết $5$ nhớ $1$.\\
-- $4 + 2 = 6$, thêm $1$ là $7$, viết $7$.
\end{minipage}

\textbf{2. Cộng ra số 100 000:}\\
$60\,000 + 40\,000 = \mathbf{100\,000}$.
}
\end{tcolorbox}
\end{minipage}
\hfill
\begin{minipage}[t]{0.49\textwidth}\vspace{0pt}
\begin{tcolorbox}[
  enhanced,
  colback=red!4!white,
  colframe=red!70!black,
  arc=2mm, boxrule=0.8pt,
  left=2mm, right=2mm, top=1.8mm, bottom=1.8mm,
  title={\footnotesize\bfseries\color{white}\faExclamationTriangle\ Mẹo tránh bẫy trắc nghiệm}
]
{\footnotesize
\begin{itemize}\setlength{\itemsep}{1.5pt}\setlength{\leftskip}{-2.5mm}
  \item \textbf{Bẫy đặt tính lệch hàng:}\\
  Khi cộng số có 5 chữ số với số có 4 chữ số (ví dụ $35\,240 + 6180$), số $6$ phải thẳng cột với số $5$ (hàng nghìn), tuyệt đối không thẳng hàng chục nghìn!
  \item Luôn cộng số nhớ vào hàng liền kề bên trái.
\end{itemize}
}
\end{tcolorbox}
\end{minipage}


\end{document}
